% Supplementary material — Extended Threats to Validity (§7.5, §7.6)
% Extracted from adaptive-deployment-control.tex for the 10-page main-paper limit.
\documentclass[10pt]{article}
\usepackage[margin=1in]{geometry}
\usepackage{amsmath,amssymb}
\providecommand{\tightlist}{\setlength{\itemsep}{0pt}\setlength{\parskip}{0pt}}
\begin{document}
\subsection{7.5 Drift Detection Conclusions Are Conditional on
PageHinkley
Calibration}\label{drift-detection-conclusions-are-conditional-on-pagehinkley-calibration}

Page-Hinkley at $\lambda$\_PH = 50 fires 44 false alarms on 1,150 stationary
steps in the ablation experiment, and 10.6 resets per 500-step
trajectory in the drift-mode evaluation. The threshold was set by the
PageHinkley default and never tuned to the cost stream's variance. All
findings about drift adaptation performance --- specifically that
\texttt{linucb\_with\_drift\_full} is 8--28\% more expensive than plain
LinUCB across all three drift modes --- are conditional on this
calibration. A lower $\lambda$\_PH (fewer false alarms) or a different detector
(ADWIN; Bifet \& Gavald{\`a}, 2007) may change the direction of the drift
finding. We report the result for $\lambda$\_PH = 50 as a measured limit of this
detector on this cost stream; we make no claim about whether drift
adaptation is intrinsically harmful.

\subsection{7.6 Thompson Sampling Propensities Are
Intractable}\label{thompson-sampling-propensities-are-intractable}

The true selection probability for Thompson Sampling requires
integrating over the posterior --- intractable in closed form. We report
propensity = 1.0 throughout. IPS correction cannot be applied to
Thompson results even in principle. Thompson's cost estimates carry the
standard replay bias plus this additional limitation.



\subsection{7.7 Seed Counts Below 30 Produce Misleading Direction
Findings}\label{seed-counts-below-30-produce-misleading-direction-findings}

Seed counts below approximately 30 produced misleading direction
findings in our pilot runs (n=5 showed Thompson as the cheapest policy
on synthetic data; n=30 reverses the direction with a non-overlapping
CI). All headline numbers in this paper use $\geq$30 seeds and bootstrap seed
42.

\subsection{7.8 CI Outcome Is a Proxy for Deployment Failure, Not a
Direct
Measure}\label{ci-outcome-is-a-proxy-for-deployment-failure-not-a-direct-measure}

The entire evaluation chain --- cost computation, policy update, and
result reporting --- treats the logged CI outcome (pass/fail) as a proxy
for the deployment failure outcome a human engineer would observe in
production. This proxy relationship is imperfect in both directions. A
CI pass does not guarantee safe deployment: integration-test gaps,
environment-specific faults, and configuration drift can produce
production incidents on passing builds. A CI failure does not guarantee
a production incident would have occurred: flaky tests, dependency
version mismatches, and transient infrastructure errors produce false CI
failures. The cost model therefore mis-prices a fraction of deployment
decisions. The direction and magnitude of this mis-pricing depend on the
test suite quality and build environment stability of each repository,
neither of which is measured here. All cost claims should be read as
costs under the CI-outcome proxy, not costs under the true
deployment-failure outcome.

\subsection{7.9 Real-Data Evaluation Assumes Stationarity; No Drift
Analysis Was
Performed}\label{real-data-evaluation-assumes-stationarity-no-drift-analysis-was-performed}

The GitHub Actions evaluation treats each 300-step project trajectory as
a stationary distribution and uses all steps for both learning and
evaluation. If the true CI failure rate changes over the 300-step window
--- for example, due to a large dependency upgrade, a project refactor,
or infrastructure migration --- then the earlier steps of the trajectory
correspond to a different distribution than the later steps. In that
case, a policy that learns from early steps may be mis-calibrated for
later steps, producing an underestimate of the cost gap between adaptive
and non-adaptive policies. We performed no stationarity test on the
real-data trajectories. The drift evaluation in §6.4 is conducted
exclusively on synthetic data; its findings do not imply anything about
whether the real GitHub Actions trajectories are stationary.

\subsection{7.10 Synthetic Failure Rates Are Unobservable Hidden
State}\label{synthetic-failure-rates-are-unobservable-hidden-state}

The synthetic dataset assigns failure rates at the project level
(smoke/alpha: 15\%, smoke/beta: 35\%), but these rates are not part of
the feature vector presented to the policy. From the policy's
perspective, the failure rate is a latent variable that must be inferred
from the reward history. This creates a structural advantage for
policies that converge quickly on the majority class: after seeing
enough failures, LinUCB correctly assigns high probability to the block
arm for smoke/beta, but the feature vector provides no direct signal
that smoke/beta is a high-failure project. In a real setting, project
metadata (team size, language, test coverage) could make the failure
rate more directly observable. The synthetic evaluation therefore
measures a lower bound on bandit performance: a policy with access to
project-type features could converge faster. The aggregate tie on the
synthetic dataset does not reflect what would happen with a richer
feature representation.

\subsection{7.11 Bootstrap Results Depend on a Single Fixed Seed;
Different Seeds May Yield Different
p-Values}\label{bootstrap-results-depend-on-a-single-fixed-seed-different-seeds-may-yield-different-p-values}

All bootstrap confidence intervals and paired bootstrap p-values in this
paper use seed 42 with 10,000 resamples (as noted in §2 and
\texttt{experiments/run\_robustness.py}). The bootstrap procedure is
deterministic given this seed. We did not evaluate sensitivity of
p-values or CI endpoints to the choice of bootstrap seed. For the F9
Thompson-vs-LinUCB result (p\textless0.0001, 0/10,000 bootstrap samples
with Thompson mean $\geq$ LinUCB), the finding is robust: the p-value is at
the floor for any 10,000-resample procedure and will not change under
different seeds. For findings where the p-value is near a threshold ---
particularly the Thompson-vs-static ``tied'' result, where static=644.5
falls inside CI {[}598, 648{]} --- a different bootstrap seed could
shift the CI endpoints by a few units and change the coverage
interpretation. Readers should treat the {[}598, 648{]} endpoints as
approximate to $\pm$5 units under seed variation.


\end{document}
